\section{Human Evaluation}
\label{apd:human-evaluation}

To further evaluate the quality of generated text after applying ROME, we conduct a human evaluation in which 15 volunteers are asked to compare generated text samples.  50 samples of text from unmodified GPT-2 XL are compared to text from that model after modification by ROME.  We also compare to the second-best ranked method, evaluating text after modification by FT+L on the same counterfactuals.  Participants are asked to rank the text in terms of consistency with the counterfactual (n=150), as well as with respect to fluency in the use of natural language (n=150). Results are summarized in \reffig{apd-human-pairwise}, and randomly-sampled examples are shown in Figures~\ref{fig:apd-human-case-1},~\ref{fig:apd-human-case-2},~\ref{fig:apd-human-case-3}.

Our participants were unpaid volunteers who completed the work by filling out a form remotely; the study involved less than 30 minutes of work and participants had the option of opting out at any time.  \reffig{apd-human-study-instructions} shows the full instructions.

\begin{figure*}
\centering
\includegraphics[width=\textwidth]{figs/appendix-human_pairwise-crop.pdf}%
\caption{\small Results from a human evaluation of generated text after applying ROME.  Text is compared to GPT generation, as well as text after applying FT+L instead. Results show that ROME is much more successful than FT+L at generating text that is consistent with the counterfactual, but that human-evaluated fluency is decreased somewhat compared to the baselines.  Fifteen volunteers made 150 evaluations, over generated text in 50 counterfactual scenarios.}
\lblfig{apd-human-pairwise}
\end{figure*}
\begin{figure*}
\includegraphics[width=\textwidth]{figs/appendix-human_case_1-crop.pdf}%
\caption{\small Human evaluation, random sample 1.}
\lblfig{apd-human-case-1}
\end{figure*}
\begin{figure*}
\includegraphics[width=\textwidth]{figs/appendix-human_case_2-crop.pdf}%
\caption{\small Human evaluation, random sample 2.}
\lblfig{apd-human-case-2}
\end{figure*}
\begin{figure*}
\includegraphics[width=\textwidth]{figs/appendix-human_case_3-crop.pdf}%
\caption{\small Human evaluation, random sample 3.}
\lblfig{apd-human-case-3}
\end{figure*}
\begin{figure*}
\includegraphics[width=0.96\textwidth]{figs/human_study_design-crop.pdf}%
\caption{\small Human evaluation, full instructions.}
\lblfig{apd-human-study-instructions}
\end{figure*}

We observe that ROME is much more successful than FT+L at generating text that is consistent with the counterfactual; this finding is consistent results in Table~\ref{tab:cf-results} that show that ROME generalizes better than FT+L.  Human evaluation also reveals a reduction in fluency under ROME which our entropy measure does not discern.  Some of the differences are subtle: examples of fluency losses detected by human raters can be seen in Figures~\ref{fig:apd-human-case-1},~\ref{fig:apd-human-case-2}.
